% Ampliación para insertar después de hibridacion.tex y del módulo térmico.
% La ecuación electrónica completa y las torres permanecen en sus secciones.
\section{Historia electrónica y conservación de la masa bajo transporte}
\label{md:r27:electron}

La ecuación \eqref{md:compos:electron-completo} evalúa una sección central
persistente y conserva su retorno de acción. La trayectoria, su memoria y
su masa son lecturas diferenciadas del mismo estado APP--TRIT--TPK. El
transporte de una ruta puede modificar eventos y memoria conservando el valor
del lector másico. Para decidirlo se utiliza el carácter completo de
\eqref{md:hib:caracter}, junto con el exponente de retorno.

\subsection{Separación y restitución de las dos historias centrales}
Los registros expresados de ambas orientaciones son
\begin{equation}
 e_{++}=(2,2,5,-4,-3,-2)^2,\qquad
 e_{--}=(4,3,2,-2,-2,-5)^2,
 \label{md:r27:historias}
\end{equation}
donde el exponente $2$ indica concatenación de dos bloques. Comparten total
nulo y el mismo histograma absoluto. Para un registro $e\in\mathbb Z^{12}$,
sean
\begin{equation}
 \begin{gathered}
 p_i=e_i+e_{i+6},\quad \mathfrak a_i=e_i-e_{i+6},\\
 s_C=\sum_{i=0}^{5}p_i,\quad M_i=p_i-p_5\quad(0\le i<5).
 \end{gathered}
 \label{md:r27:registro-electron}
\end{equation}
\begin{proposition}[Restitución integral del registro]
El dato $(s_C,M_0,\ldots,M_4,\mathfrak a_0,\ldots,\mathfrak a_5)$ determina el registro mediante
\begin{equation}
 \begin{gathered}
 p_5=\frac{s_C-\sum_{i=0}^{4}M_i}{6},\quad p_i=p_5+M_i,\\
 e_i=\frac{p_i+\mathfrak a_i}{2},\quad e_{i+6}=\frac{p_i-\mathfrak a_i}{2}.
 \end{gathered}
 \label{md:r27:inversa-electron}
\end{equation}
Su imagen integral se caracteriza por la divisibilidad por seis del primer
numerador y la congruencia $p_i\equiv \mathfrak a_i\pmod2$ para cada $i$.
\end{proposition}
\begin{proof}
Sumar $p_i=p_5+M_i$ para los cinco primeros índices y añadir $p_5$ da
$s_C=6p_5+\sum M_i$. Sumar y restar las dos definiciones de $p_i,\mathfrak a_i$
da las fórmulas restantes. Las condiciones indicadas son necesarias para
obtener enteros y suficientes para reconstruirlos; la sustitución devuelve
todos los datos originales.
\end{proof}
Para \eqref{md:r27:historias}, $\mathfrak a=0$ y
$M_{++}=(8,8,14,-4,-2)$, mientras $M_{--}=(18,16,14,6,6)$.
La restitución distingue ambas historias aunque lectores escalares concretos
compartan su valor. Los símbolos $\mathfrak a_i$ de esta descomposición son diferencias
de eventos y se distinguen del carácter areal $a_0$ por su dominio.

\subsection{Forma espectral del lector de retorno}
Para $\tau\in\mathbb R^{12}$, con índices cíclicos, definimos
\begin{equation}
 C_s=\sum_{j=0}^{11}\tau_j\tau_{j+s},\quad
 T_k=\sum_{j=0}^{11}\tau_j e^{-2\pi i kj/12},\quad
 \theta_k=2\pi k/12.
 \label{md:r27:fourier}
\end{equation}
La representación por caracteres se aplica después de producir el registro.
El lector de retorno conserva
$\Omega=-2C_1-4C_2-6C_3$ y $P_6=C_6/2$.
\begin{proposition}[Descomposición armónica exacta]
\begin{equation}
 \begin{split}
 C_s&=\frac1{12}\sum_{k=0}^{11}|T_k|^2\cos(s\theta_k),\\
 \Omega&=\frac1{12}\sum_{k=0}^{11}|T_k|^2
 [-2\cos\theta_k-4\cos2\theta_k-6\cos3\theta_k],\\
 P_6&=\frac1{24}\sum_{k=0}^{11}|T_k|^2(-1)^k.
 \end{split}
 \label{md:r27:espectro}
\end{equation}
\end{proposition}
\begin{proof}
Desarrollar $|T_k|^2$ produce
$\sum_{j,l}\tau_j\tau_l e^{-i\theta_k(j-l)}$.
La suma $\sum_{k=0}^{11}e^{i\theta_km}$ vale doce si $12$ divide $m$ y
cero en otro caso. Aplicar esta identidad al desplazamiento $s$ recupera
$C_s$. Como $\tau$ es real, las partes imaginarias se cancelan en pares
$k,12-k$ y queda el coseno. Las otras dos fórmulas resultan por combinación
lineal y por $\cos(6\theta_k)=(-1)^k$.
\end{proof}
Para $\tau_e=(+++---)^2$, sólo quedan las potencias
$|T_2|^2=64$, $|T_6|^2=16$, $|T_{10}|^2=64$. Sus pesos en $\Omega$
son $7,4,7$, de modo que $\Omega=80$ y $P_6=6$.
Se recupera así la parte armónica de
$\kappa_e=80-54\alpha+6\Delta_4$ en la ecuación electrónica íntegra.
El registro alternante $(+-)^6$, con igual histograma, tiene
$|T_6|^2=144$ y $\Omega=48$: el orden es una dependencia efectiva del lector.
Las traslaciones cíclicas conservan $|T_k|^2$, aunque cambien la fase de
$T_k$; esta lectura espectral preserva correlaciones, no toda la historia.

\subsection{Criterio de masa transportada y sus lecturas asociadas}
Sea $\mathcal H_{\mathcal F}$ la fibra de un dominio finito de rutas y
$\{|\Gamma\rangle\}$ su base. Se fija uno de los lectores de retorno
$r=D$ o $r=\Omega$ y una coordenatización común con $R_{\rm act}>0$.
La ley de masas de \eqref{md:hib:factor-ruta} determina
\begin{equation}
 \begin{gathered}
 \frac{m_\Gamma^\beta}{m_e^\beta}
 =\chi_{\rm ref}(\sigma_\Gamma-\sigma_e,\eta_\Gamma-\eta_e)
 R_{\rm act}^{\kappa_\Gamma-\kappa_e},\\
 \Lambda_\Gamma=\ln(m_\Gamma^\beta/m_e^\beta).
 \end{gathered}
 \label{md:r27:masa-relativa}
\end{equation}
El logaritmo actúa sobre una razón adimensional. En una fibra no escalar
se mantiene el operador y su realización diagonal declarada; el dominio
finito puede ser una órbita o multisección, sin seleccionar una celda arbitraria.

\begin{theorem}[Conservación del lector másico en un dominio estable]
Sea $F$ un transporte TPK que actúa biyectivamente en el dominio fijado y
$V|\Gamma\rangle=|F\Gamma\rangle$. Para
$M|\Gamma\rangle=m_e^\beta e^{\Lambda_\Gamma}|\Gamma\rangle$, se cumple
\begin{equation}
 [M,V]|\Gamma\rangle=m_e^\beta
 (e^{\Lambda_{F\Gamma}}-e^{\Lambda_\Gamma})|F\Gamma\rangle.
 \label{md:r27:conmutador-masa}
\end{equation}
Por tanto, $[M,V]=0$ equivale a $\Lambda_{F\Gamma}=\Lambda_\Gamma$ para
todas las rutas. Con los parámetros de la coordenatización fijos, la condición es
\begin{equation}
 \begin{aligned}
 &\delta n_A x+\frac{\delta s_C}{6}y+
 \delta k_{\Delta_4}\Delta_4+\delta b_\zeta\zeta_{270}\\
 &\qquad+\sum_h\delta N_h(q_-^h-q_+^h)
 +\delta\kappa\ln R_{\rm act}=0,
 \end{aligned}
 \label{md:r27:criterio-masa}
\end{equation}
donde $\delta$ compara la ruta transportada con la inicial.
\end{theorem}
\begin{proof}
Aplicar $MV$ y $VM$ a cada vector de la base da
\eqref{md:r27:conmutador-masa}. La positividad de $m_e^\beta$ y la
inyectividad de la exponencial real dan la equivalencia. Tomar el logaritmo
de \eqref{md:r27:masa-relativa} y restar los caracteres de ambas rutas da
\eqref{md:r27:criterio-masa}. En torres infinitas se exige para cada ruta
$\sum_h|N_h(q_-^h-q_+^h)|<\infty$; esta convergencia absoluta legitima
la resta. La finitud del dominio de rutas conserva los operadores acotados
utilizados en esta prueba.
\end{proof}

El transporte se produce primero mediante TPK; la igualdad decide después
qué conserva. Los eventos individuales y la memoria pueden cambiar y
compensarse en \eqref{md:r27:criterio-masa}. Al mantener la sección de acción,
el reloj y la realización circular de \eqref{md:hib:reciprocidad-radial},
la conmutación con $M$ conserva también las lecturas de frecuencia propia,
radio gravitatorio, longitud de Compton reducida y peso térmico que son
funciones del mismo operador positivo. Para el inverso de masa se trabaja
fuera del sector nulo. En la realización marcada de
\eqref{md:r27:area-gravedad}, la cuarta lectura adopta la forma
\begin{equation}
 \Xi_\Gamma=\frac{\Omega_\Gamma}{\omega_0}
 =\frac{r_{g,\Gamma}}{L_\alpha}
 =\frac{L_\alpha}{\bar\lambda_\Gamma}
 =\frac{E_\Gamma}{b_*\Theta_{{\rm clk},P}\ln3}.
 \label{md:r27:lecturas-conservadas}
\end{equation}
La prueba es sustitución de las igualdades ya demostradas en una misma coordenatización.
Este criterio describe conservación durante un transporte declarado; los
selectores físicos de especie, de representación y de interacción conservan
sus condiciones propias. La trayectoria y su lectura armónica añaden
información a la masa escalar sin reemplazar esas condiciones.
